\chapter{Related Work}
\label{ch:related}

Automated brain tumour segmentation has been studied for over two decades, and
the field has moved from hand-engineered pipelines to end-to-end deep learning.
This chapter surveys the methods most relevant to the present work and situates
the chosen approach among them.

\section{Classical and early machine-learning methods}

Early automated approaches relied on intensity thresholding, region growing,
atlas registration and clustering (for example $k$-means or fuzzy $c$-means),
often combined with hand-crafted texture or intensity features fed to classifiers
such as support vector machines or random forests. These methods are
interpretable and require no large training set, but they are sensitive to
intensity inhomogeneity and noise, depend heavily on manual feature design and
parameter tuning, and generally fail to capture the strong spatial context needed
to separate diffuse tumour boundaries from healthy tissue. The Multimodal Brain
Tumour Segmentation (BraTS) benchmark was introduced in part to compare such
methods on a common footing and quantified their limitations against expert
annotation \citep{menze2015brats}.

\section{Fully convolutional and encoder--decoder networks}

The decisive shift came with fully convolutional networks and, in particular, the
U-Net \citep{ronneberger2015unet}. The U-Net couples a contracting encoder that
captures context with a symmetric expanding decoder that restores spatial
resolution, and it connects the two with skip connections that carry
high-resolution features across the bottleneck. This design is highly effective
for biomedical segmentation, where training data are scarce and precise
boundaries matter.

Because brain tumours are inherently volumetric, purely 2D models that process
slices independently discard information along the through-plane axis.
Two influential extensions address this directly: the 3D~U-Net, which replaces
all 2D operations with their volumetric counterparts and learns from sparsely
annotated volumes \citep{cicek20163dunet}, and V-Net, which introduced a
Dice-based objective together with residual connections for volumetric medical
segmentation \citep{milletari2016vnet}. The present work adopts a 3D~U-Net for
exactly this reason.

A major practical lesson from the field is that careful configuration often
matters more than architectural novelty. The nnU-Net framework
\citep{isensee2021nnunet} showed that a well-tuned, self-configuring U-Net --- with
principled choices for pre-processing, patch size, normalisation and training ---
matches or exceeds far more elaborate models across many medical segmentation
tasks, and it remains a strong baseline on BraTS. This motivates our emphasis on
a solid, standard 3D~U-Net with sound pre-processing and training rather than a
bespoke architecture.

\section{Attention and transformer-based models}

More recent work introduces attention to focus computation on salient regions and
to model long-range dependencies. Attention U-Net adds gating on the skip
connections so that irrelevant background features are suppressed
\citep{oktay2018attentionunet}. Transformer-based hybrids bring global
self-attention into the encoder: TransBTS couples a convolutional encoder with a
transformer bottleneck for 3D brain tumour segmentation \citep{wang2021transbts},
and Swin~UNETR uses a hierarchical shifted-window transformer encoder with a
convolutional decoder and reports strong BraTS performance
\citep{hatamizadeh2022swinunetr}. These models can improve accuracy but are more
data- and compute-hungry; on a single consumer GPU with limited memory a compact
3D~U-Net is a more practical starting point, with transformer encoders left as a
natural extension.

\section{The BraTS benchmark}

The BraTS challenge has been the central benchmark for this task
\citep{menze2015brats,bakas2017advancing}. It provides multi-institutional,
multi-parametric MRI scans that are skull-stripped, co-registered to a common
anatomical template and resampled to $1\,\mathrm{mm}^3$ isotropic resolution, with
expert annotations of three tumour sub-regions. Crucially for this project, the
BraTS~2020 release labels each case by grade, distinguishing high-grade glioma
(glioblastoma) from lower-grade glioma; this label makes it possible to restrict
training and evaluation to genuine glioblastoma cases. Standardised data and
evaluation make results across methods directly comparable, and the Dice
coefficient and 95th-percentile Hausdorff distance used here are the benchmark's
established metrics.

\section{Positioning of this work}

Table~\ref{tab:related} summarises the main architectural families. This project
deliberately builds on a standard 3D~U-Net rather than a transformer model: the
goal is a correct, reproducible and honestly evaluated glioblastoma-specific
pipeline that runs on a single 8~GB GPU, with clear headroom to adopt the more
advanced encoders discussed above as future work.

\begin{table}[t]
\centering
\caption{Representative architectural families for brain tumour segmentation.}
\label{tab:related}
\small
\setlength{\tabcolsep}{4pt}
\begin{tabularx}{\textwidth}{@{}l R R@{}}
\toprule
Family & Representative work & Key idea \\
\midrule
2D encoder--decoder & U-Net \citep{ronneberger2015unet} & Skip connections, per-slice \\
3D encoder--decoder & 3D~U-Net \citep{cicek20163dunet}, V-Net \citep{milletari2016vnet} & Volumetric context, Dice loss \\
Self-configuring & nnU-Net \citep{isensee2021nnunet} & Automated pipeline configuration \\
Attention-gated & Attention U-Net \citep{oktay2018attentionunet} & Gated skip connections \\
Transformer hybrid & TransBTS \citep{wang2021transbts}, Swin~UNETR \citep{hatamizadeh2022swinunetr} & Global self-attention encoder \\
\bottomrule
\end{tabularx}
\end{table}
